\documentclass[conference]{IEEEtran}
\IEEEoverridecommandlockouts
% The preceding line is only needed to identify funding in the first footnote. If that is unneeded, please comment it out.
\usepackage{cite}
\usepackage{amsmath,amssymb,amsfonts}
\usepackage{algorithmic}
\usepackage{graphicx}
\usepackage{textcomp}
\usepackage{xcolor}
\usepackage{booktabs}
\usepackage{cleveref}
\def\BibTeX{{\rm B\kern-.05em{\sc i\kern-.025em b}\kern-.08em
    T\kern-.1667em\lower.7ex\hbox{E}\kern-.125emX}}
\begin{document}

\title{Paper Title\\
\thanks{This thesis was prepared in partial fulfilment of the requirements for the Degree of Bachelor of Computer Science, Maastricht University. Supervisor: Petra Bosilj}
}

\author{\IEEEauthorblockN{Adriane Bien}
\IEEEauthorblockA{\textit{Department of Advanced Computing Sciences} \\
\textit{Faculty of Science and Engineering}\\
\textit{Maastricht University}\\
Maastricht, The Netherlands}
}

\maketitle

\begin{abstract}
This document is a model and instructions for \LaTeX.
This and the IEEEtran.cls file define the components of your paper [title, text, heads, etc.]. *CRITICAL: Do Not Use Symbols, Special Characters, Footnotes, 
or Math in Paper Title or Abstract.
\end{abstract}

\begin{IEEEkeywords}
component, formatting, style, styling, insert
\end{IEEEkeywords}

\section{Introduction}
\label{sec:intro}
To meet the growing global demand for protein, insect farming has emerged as a
promising alternative to conventional livestock production~\cite{vanhuis2013}.
\textit{Tenebrio molitor}, the yellow mealworm, is one of the few insect species
approved for human consumption in Europe~\cite{eu2021tenebrio}. To be
commercially viable, however, production must scale significantly---a challenge
that hinges directly on automation~\cite{rumpold2013}.

Automated mealworm farming requires continuous monitoring of the rearing plates
to inform decisions about feeding, harvesting, and environmental control. Prior
work has explored machine vision for this purpose, but most of it relies on
object detection and instance segmentation~\cite{majewski2022multipurpose,
Papadopoulos2024TenebrioVision}, which notoriously struggle with the dense,
heavily occluded scenes typical of insect rearing trays. Crowd counting via
density estimation offers a far more suitable framework for these
conditions~\cite{Sindagi2017ASO}, yet it remains entirely unexplored in insect
agriculture.

The primary driver for automation is economic: reducing operational costs and
maximizing production efficiency~\cite{rumpold2013}. Because the industry lacks
standardization~\cite{NAWOYA2024108503}, farms have adopted a diverse range of
technological solutions~\cite{CoRoSect2022AdvancedArchitecture,
kroncke2020automation}. One compelling approach is edge AI, which is
particularly attractive for cost-sensitive operations because it runs inference
locally on low-cost embedded devices~\cite{angulo2025realtimecricketsortingsex}.
This thesis therefore investigates how crowd-counting models can be deployed
within the strict memory constraints of edge hardware.

Memory consumption during inference is governed primarily by two factors: model
complexity and input resolution. Larger architectures require more memory to
store their parameters, while higher-resolution inputs enlarge the intermediate
feature maps that must be held in memory during the forward
pass~\cite{Chao2019HarDNet}. Reducing the memory footprint therefore comes down
to two options: shrinking the model or shrinking the input.

This raises a practical question: when memory is strictly limited, is
architectural complexity or input fidelity more important? To investigate, this
thesis compares two competing strategies---deploying a ``smart but heavy'' model
that must be fed downscaled images to fit within memory, against a ``simpler but
light'' model that leaves enough headroom to process images at full resolution.

\section{Related Works}
\subsection{Crowd Counting Approaches}
Crowd counting methods can be broadly categorized into three approaches:
detection-based, regression-based, and density estimation-based methods.

The earliest approaches were detection-based, producing a population estimate by
aggregating individual bounding-box predictions. As noted above, these methods
struggle considerably under severe occlusion and high object density. %can mention specific failures from the papers.

Regression-based methods were developed to address this by computing object
counts directly from image features, bypassing the need for explicit
localization. While more robust to occlusion, they sacrifice spatial information
entirely, giving no insight into how individuals are distributed across the
scene.

The most widely used crowd-counting models today are based on density
estimation. Rather than localizing individuals explicitly, this approach learns
to map an image to a density map, where the integral over any region gives the
count within it and the integral over the entire map gives the total count.
Because the map preserves where mass is concentrated across the scene, density
estimation retains the spatial information that regression discards. Models are
trained by regressing the predicted map against a ground-truth density
map~\cite{wang2020mobilecount, li2018csrnet, gao2020crowdcountingsurvey}.

\subsection{CNN-Based Density Estimation}
The success of Convolutional Neural Networks (CNNs) in image classification and
recognition has made them the de facto standard for density estimation. Rather
than relying on hand-crafted features and shallow regressors, CNNs learn
hierarchical feature representations directly from raw pixel data, enabling far
greater robustness to occlusion, lighting variation, and background
clutter~\cite{gao2020crowdcountingsurvey, li2018csrnet}.

Early approaches used basic CNNs with standard convolutional and fully connected
layers. These were simple to implement but struggled with the scale variation
inherent to crowd scenes, motivating multi-column architectures such as MCNN,
which run parallel branches with different receptive fields to capture features
at multiple scales~\cite{Zhang_2016_CVPR}.

The field has since largely shifted toward single-column designs, with CSRNet a
pivotal model in this transition. CSRNet demonstrated that a single deeper column
could outperform bloated multi-column structures through a dilated-convolution
back-end that enlarges the receptive field while preserving spatial resolution.
This design has since become one of the most widely adopted in crowd counting,
and CSRNet remains a standard baseline for benchmarking newer methods,
particularly in dense scenes~\cite{gao2020crowdcountingsurvey, li2018csrnet}.

\subsection{Lightweight Networks}
Beyond architectural design, lightweight models have emerged as a crucial area of
research for resource-constrained deployment. The objective is to minimize
computational overhead and parameter count without incurring a significant
penalty to counting accuracy~\cite{gao2020crowdcountingsurvey,
wang2020mobilecount}.

A prominent example in density-based crowd counting is MobileCount, an
encoder-decoder network built on a lightweight MobileNetV2 backbone. Compared
with CSRNet, MobileCount reduces the model footprint by roughly 80\%---only
$3.40$ million parameters versus CSRNet's $16.26$ million. As shown in
Table~\ref{tab:lit_table}, MobileCount exhibits only a minor reduction in
accuracy across standard benchmark datasets, remaining highly comparable to its
heavier counterpart. Given that it requires 52 times fewer floating-point
operations (FLOPs), this marginal trade-off in accuracy is heavily outweighed by
its gains in computational efficiency~\cite{wang2020mobilecount}.

\begin{table}[t]
    \centering
    \caption{Reported MAE for CSRNet and MobileCount on standard
    crowd-counting benchmarks. SHT: ShanghaiTech; WE: WorldExpo'10.
    ~\cite{wang2020mobilecount, li2018csrnet}}
    \label{tab:lit_table}
    \footnotesize
    \setlength{\tabcolsep}{5pt}
    \begin{tabular}{lcccc}
        \toprule
        Network     & SHT-A & SHT-B & UCF-CC-50 & WE \\
        \midrule
        CSRNet      & 68.2 & 10.6 & 266.1 & 8.6  \\
        MobileCount & 89.4 & 9.0  & 284.8 & 11.1 \\
        \bottomrule
    \end{tabular}
\end{table}

\section{Methods}
\subsection{Dataset}
This study uses the TenebrioVision dataset, developed within the CoRoSect
project. It consists of 1{,}120 high-resolution images ($3088 \times 2076$
pixels), grouped into 14 count levels of 80 images each, where every image in a
level contains a fixed number of \textit{Tenebrio molitor} larvae: 10, 15, 20,
25, 30, 35, 40, 45, 50, 60, 70, 80, 90, or 100. Each larva is labelled with a
bounding box, giving 53{,}600 annotated instances in total. TenebrioVision is
currently the only publicly available dataset reflecting the high-density
conditions typical of industrial larvae farming, making it the most suitable
benchmark for crowd-counting models in this specific
domain~\cite{Papadopoulos2024TenebrioVision}.

Following the split protocol of~\cite{Papadopoulos2024TenebrioVision}, the data
is stratified by count level: within each level the 80 images are shuffled with
a fixed random seed (42) and partitioned into 60 training, 12 validation, and 8
test images. Repeating this across all 14 levels yields 840 training, 168
validation, and 112 test images---a $75\,/\,15\,/\,10\,\%$ split in which every
count level is represented in identical proportion across all three subsets.
This stratification ensures a balanced distribution, so differences in
validation or test error cannot be attributed to one subset containing
systematically denser or sparser images.

\subsection{Models}
Two CNN-based density-estimation models are trained and evaluated: the accurate but computationally heavy CSRNet, and the lightweight, efficiency-oriented MobileCount. Together they represent complementary points on the accuracy–efficiency trade-off.

\subsubsection{CSRNet}
CSRNet's architecture consists of a front-end feature extractor and a dilated
convolutional back-end. The front-end is built from the first ten layers of a
pre-trained VGG-16 network, with the fully connected classification layers
removed. Three max-pooling operations reduce the spatial resolution to
$H/8 \times W/8$, producing a feature map that benefits from strong transfer
learning capacity while retaining meaningful spatial structure.

Rather than continuing to downsample via pooling---which would further degrade
spatial resolution and make accurate density-map reconstruction
difficult---the back-end uses dilated convolutions to expand the receptive field
without reducing resolution or adding parameters. A dilated convolution with
dilation rate $r$ effectively enlarges a $k \times k$ kernel to
$k + (k-1)(r-1)$, allowing the network to aggregate multi-scale contextual
information while keeping feature maps at a constant spatial size. CSRNet is proposed in four configurations with varying back-end dilation rates; configuration~B is used here, which the original paper identifies as most accurate and which alternates between dilation rates of 2 and 4 throughout the back-end layers.

A final $1 \times 1$ convolution collapses the back-end feature maps to a single
channel, producing a density map at $H/8 \times W/8$~\cite{li2018csrnet}.

\subsubsection{MobileCount}
MobileCount's encoder is based on MobileNetV2, a backbone built around an
inverted residual structure with linear bottlenecks. Each bottleneck applies a
$1 \times 1$ pointwise convolution to expand channel dimensionality, a
$3 \times 3$ depthwise convolution to capture spatial features, and a second
$1 \times 1$ pointwise convolution to project back to a lower-dimensional
representation. Separating channel mixing from spatial filtering in this way
dramatically reduces parameter count relative to standard convolutions.

To reduce computation further, MobileCount uses only four of MobileNetV2's seven
bottleneck blocks, as the original paper reports that removing the final three
improves performance alongside a substantial reduction in FLOPs. An additional
$3 \times 3$ max-pooling layer of stride 2 is inserted at the start of the
encoder, reducing input resolution early in the pipeline and cutting FLOPs by
over 70\% at marginal accuracy cost. Together with the initial stride-2
convolution, this reduces the resolution at the first bottleneck block to
$H/4 \times W/4$; subsequent blocks downsample further to $H/8$, $H/16$, and
$H/32$.

The decoder is a simplified LightWeight RefineNet, originally developed for
real-time semantic segmentation. It progressively fuses the four multi-scale
encoder feature maps through a series of Chained Residual Pooling (CRP) and
fusion blocks, upsampling from $H/32$ back through each scale until reaching
$H/4$. A final $1 \times 1$ convolutional layer maps this representation to a
single-channel output at $H/4 \times W/4$, which is then bilinearly upsampled
$4\times$ to full input resolution to produce the predicted density
map~\cite{wang2020mobilecount}.

\subsection{Loss Function}

Both models are trained end-to-end using the mean squared error loss between the predicted density map $D^{PM}$ and the ground truth density map $D^{GT}$:

\begin{equation}
    \mathcal{L} = \frac{1}{N} \sum_{i=1}^{N} \| D^{PM}(i) - D^{GT}(i) \|_2^2
\end{equation}

where $N$ is the number of training samples. The final population estimate is obtained by integrating the predicted density map over the full image. \cite{wang2020mobilecount, li2018csrnet}

\subsection{Image Resolution}
To investigate the impact of input scale on counting performance, both models are evaluated across a hierarchy of six resolutions. These are derived by progressively halving the original $3088 \times 2076$ images (rounding up on odd dimensions), giving: $1544 \times 1038$, $772 \times 519$, $386 \times 260$, $193 \times 130$, $97 \times 65$, and $49 \times 33$.

\subsection{Ground Truth Generation}
Point annotations are derived from the centre of each bounding box to construct
a discrete point map. Ground-truth density maps are then generated by convolving
this point map with a Gaussian kernel. Formally, for a set of $N$ annotated
positions, the density map $D(\mathbf{x})$ is defined as
\begin{equation}
    D(\mathbf{x}) = \sum_{i=1}^{N} \delta(\mathbf{x} - \mathbf{x}_i) *
    G_{\sigma}(\mathbf{x}),
\end{equation}
where $\delta(\mathbf{x} - \mathbf{x}_i)$ is a Dirac delta centred at larva
position $\mathbf{x}_i$ and $G_{\sigma}$ is an isotropic Gaussian kernel with
standard deviation $\sigma$.

A fixed-sigma kernel was applied to every annotation, following the strategy
CSRNet uses on ShanghaiTech Part~B. That dataset's structural
characteristics---a fixed viewpoint, roughly constant object size, and
comparatively uniform density---closely mirror those of TenebrioVision, and
CSRNet models it with a fixed kernel of $\sigma = 15$, approximately half the
typical head diameter. The same object-to-kernel ratio is adopted here: the standard
deviation is set to $\sigma = \texttt{bbox}/2$, where \texttt{bbox} is the
average larva bounding-box size in TenebrioVision, so that $\sigma$ again spans
roughly half the object's diameter.

The kernel was truncated to an effective window of $8\sigma + 1$ pixels. To
preserve the relative spatial extent of the density distribution across scales,
$\sigma$ was scaled proportionally with resolution, giving values of 24, 12, 6,
3, 2, and 1 pixels for the six variants respectively (rounded to the nearest
pixel at the two smallest scales).

\subsection{Evaluation Metrics}
\label{sec:eval_metrics}
Following standard practice in the crowd-counting literature~\cite{gao2020crowdcountingsurvey}, Mean Absolute Error (MAE) and Mean Squared Error (MSE) are used as the primary accuracy metrics:

\begin{equation}
    \text{MAE} = \frac{1}{N} \sum_{i=1}^{N} \left| y_i - \hat{y}_i \right|,
\end{equation}
\begin{equation}
    \text{MSE} = \sqrt{\frac{1}{N} \sum_{i=1}^{N}
    \left( y_i - \hat{y}_i \right)^2},
\end{equation}
where $N$ is the number of test images, $y_i$ is the ground-truth count, and
$\hat{y}_i$ is the predicted count. MAE reflects the average magnitude of
prediction error, while MSE penalises larger errors more heavily, making it
sensitive to outliers.

\paragraph{Memory measurement.}
For each of the twelve network~$\times$~resolution cells, peak
inference memory was measured on the nvidia. The
measurement captures weights, input, and peak activations
together. Reported values are allocated peaks; the allocator's
reserved pool is somewhat higher. Because the memory footprint of a
convolutional forward pass depends only on the input dimensions, not
its content, the numbers hold for any input at a given resolution. % what does this mean?

\subsection{Experimental Design}
A model's hyperparameter configuration has a substantial effect on its final
accuracy. A fair comparison between architectures must therefore evaluate each in
a setting appropriate to it, rather than imposing a single shared configuration
that may favour one design over the other.

The original CSRNet implementation was optimised with Stochastic Gradient Descent
(SGD), a fixed learning rate of $1\times10^{-6}$, a weight decay of
$1\times10^{-4}$, and a batch size of $1$~\cite{li2018csrnet}. Subsequent work,
however, reports that the Adam optimiser yields faster convergence and stronger
final performance for this architecture even when all other original parameters
are held constant~\cite{jimaging7020021}. The original MobileCount implementation
used Adam with an initial learning rate of $1\times10^{-4}$, an exponential
learning-rate decay of $0.995$ per epoch, a weight decay of $1\times10^{-4}$, and
a batch size of $6$~\cite{wang2020mobilecount}.

All experiments are implemented in the C\textsuperscript{3}
Framework~\cite{gao2019c}, a widely used crowd-counting codebase whose authors
have already tuned the supported architectures on the standard benchmark
datasets; these configurations define the starting point for tuning
(Table~\ref{tab:baseline_hyperparameters}). Because the official MobileCount
implementation is built on top of C\textsuperscript{3}, the framework's
MobileCount configuration coincides exactly with the parameters reported in the
original publication. For CSRNet the framework instead provides a modernised
baseline; its configuration for the ShanghaiTech~Part~B dataset was adopted here.
Three axes were then tuned: the optimiser, the base learning rate, and the batch
size.

\begin{table}[t]
    \centering
    \caption{Baseline (pre-tuning) optimisation hyperparameters.}
    \label{tab:baseline_hyperparameters}
    \begin{tabular}{lcc}
        \toprule
        Hyperparameter & CSRNet & MobileCount \\
        \midrule
        Optimiser             & \multicolumn{2}{c}{Adam} \\
        Weight decay          & \multicolumn{2}{c}{$1\times10^{-4}$} \\
        Batch size            & \multicolumn{2}{c}{$6$} \\
        Initial learning rate & $1\times10^{-5}$ & $1\times10^{-4}$ \\
        LR decay (per epoch)  & \multicolumn{2}{c}{step $\gamma=0.995$} \\
        \bottomrule
    \end{tabular}
\end{table}


\subsubsection{Fixed Design Decisions}
Before any hyperparameter was swept, several decisions were frozen so that every configuration was compared on equal terms: the weight decay, the training budget and learning-rate decay, the density-target scaling, and the choice of tuning resolution. 

\paragraph{Weight decay.}
Both architectures use the \CCC{} framework's default weight decay of $1\times10^{-4}$ (\cref{tab:baseline_hyperparameters}). No preliminary
testing suggested this value was problematic, and since weight decay interacts with both the learning rate and the optimiser, tuning it jointly with them would have made the sweep infeasibly large.

\paragraph{Training budget and learning-rate decay.}
Both architectures use a fixed budget of $1500$ epochs with no early stopping, giving every configuration equal compute. With the framework's step decay of $\gamma=0.995$ per epoch, the learning rate is reduced to $0.995^{1500}\approx0.05\%$ of its base value by the end
of training. Because the decay is relative, the same epoch count is valid for every base learning rate.

\paragraph{Density-target scaling.}
The framework multiplies the density-map regression target by a
constant scaling factor $k$ (\texttt{LOG\_PARA} in the framework's
terminology) and divides the prediction by $k$ at evaluation time, so
that the per-pixel targets are of a magnitude the network can regress
comfortably. The values validated in the original works were taken as the starting point:
$k=100$ for CSRNet and $k=2550$ for MobileCount, both established on
ShanghaiTech~Part~B. Since a density map sums to the object count, its
per-pixel magnitude is inversely proportional to the number of pixels;
a value of $k$ tuned at one resolution is therefore wrong at every
other. The original values were consequently anchored to the resolution
variant closest in pixel count to ShanghaiTech~Part~B, namely
$772\times519$, and scaled in proportion to pixel count for all other
variants,
\begin{equation}
    k(r) \;=\; k_{\mathrm{base}}\cdot
    \frac{A(r)}{A(772\times519)},
    \label{eq:logpara_scaling}
\end{equation}
where $A(\cdot)$ denotes the number of pixels. 

\paragraph{Choice of tuning resolution.}
All tuning was performed on the $386\times260$ variant, the mid-range point of the six resolution variants spanning $49\times33$ to $1544\times1038$. The smaller the variant, the more image detail is lost and the less meaningful the resulting hyperparameter signal; the larger the variant, the slower the sweep (an estimated ${\approx}26$\,h per CSRNet run at the top resolution). The $386\times260$ variant is thus the cheapest \emph{healthy} regime: a full training run takes
${\approx}1$\,h for MobileCount and ${\approx}2.7$\,h for CSRNet, so dozens of configurations can be swept. Being mid-range, the selected hyperparameters must extrapolate in both directions; they are frozen after tuning and reused across all six resolutions to permit a fair cross-resolution comparison.

\subsubsection{Hyperparameter Tuning Protocol}
\label{sec:tuning_protocol}
Four axes were tuned: the augmentation scheme, the optimiser, the base
learning rate, and the batch size. The augmentation scheme was settled
first and frozen for all subsequent experiments. Optimiser and
learning rate were then swept jointly as a full grid, as preliminary
testing showed the two to be too tightly coupled to tune
independently. Batch size was swept last, around the resulting
optimum, with the learning rate scaled by the standard $\sqrt{B/6}$
rule rather than as a full grid, which would have been prohibitively
expensive. All runs use a fixed random seed ($3035$). Throughout,
validation MAE is computed after every epoch and the checkpoint with
the lowest validation MAE is selected. As a sanity criterion for the
fixed $1500$-epoch budget, a run whose best checkpoint lands in the
final ${\sim}5\%$ of epochs would indicate an insufficient budget;
this did not occur for any competitive configuration.

\paragraph{Augmentation.}
The training set contains only $840$ fixed images — small enough that
training on the full images invites memorisation, making augmentation
a necessity rather than an option. Patch-based augmentation is the
standard remedy in crowd counting, and the two reference
implementations each prescribe one: MobileCount uses online random
crops covering $80\%$ of the image, generated afresh at every
iteration~\cite{wang2020mobilecount}, while CSRNet uses an offline
procedure that fixes $18$ patches per image in
advance~\cite{li2018csrnet}.
Both schemes were compared at $386\times260$ under the baseline
hyperparameters (Table~\ref{tab:aug_paper_schemes}). Neither shows a
train--validation gap and their validation optima differ by at most
$0.03$ MAE — too little to justify a preference on accuracy grounds. 
The online crop was therefore adopted on cost and simplicity %say also its in general th way to do it now, the csrnet way is outdated because now it generates a stream that never repeats a sample:
it requires no offline expansion of the dataset, % talk about cost here - find it the time it takes for both


The crop is defined relative to the image, it
transfers unchanged across all resolution variants. Validation and test images are always processed whole. 

\begin{table}[t]
    \centering
    \caption{Augmentation schemes at $386\times260$ under the baseline
    hyperparameters. }
    \label{tab:aug_paper_schemes}
    \small
    \begin{tabular}{l rrrrrr}
        \toprule
        & \multicolumn{3}{c}{CSRNet}& \multicolumn{3}{c}{MobileCount}\\
        
        Scheme& train & val & $\Delta$ & train & val & $\Delta$ \\
        \midrule
        $80\%$ crop& $3.465$ & $3.513$ & $+0.05$ & $0.621$ & $0.632$ & $+0.011$ \\
        $18$-patch& $4.085$ & $4.074$ & $-0.01$ & $0.640$ & $0.622$ & $-0.018$ \\
        \bottomrule
    \end{tabular}
\end{table}
%this is not proof of overfitting at all.
%overfitting is when is the final MAE is way higher than the best MAE, and the training MAE gap closes, aka the training MAE keeps going down, while the val rises. 


\paragraph{Optimiser and learning rate.}
Table~\ref{tab:lr_grid} and
Figures~\ref{fig:lrxopt_csrnet}--\ref{fig:lrxopt_mobilecount} present
the joint optimiser~$\times$~learning-rate grid.
The compared optimisers, Adam and AdamW, differ only in how the
(fixed) weight decay enters the update — Adam folds it into the
gradient before the adaptive normalisation, whereas AdamW applies it
as a separate, decoupled term~\cite{loshchilov2019decoupled}.
% TODO: add loshchilov2019decoupled (AdamW paper) to your .bib.
The learning rate was swept in decades spanning one decade below and
one above each network's framework optimum
(Table~\ref{tab:baseline_hyperparameters}). Because the best cell
landed on the upper edge for both networks, the sweep was extended by
one further decade ($10^{-3}$ for CSRNet, $10^{-2}$ for MobileCount).

Both extension runs degenerated: validation MAE plateaued around 22.88 and 23.14 respectively. Training collapsed to a constant-output solution that sits right at the trivial predict-the-mean floor of 23.27. These runs are omitted from Table~\ref{tab:lr_grid}.

Two regularities emerge from the grid. First, each architecture has
its own optimal learning rate, exactly one decade apart: $10^{-4}$ for
CSRNet and $10^{-3}$ for MobileCount, in both cases one decade above
the framework default.

Second, \emph{at its optimal learning rate, each network prefers
AdamW}, but by margins that differ by an order of magnitude: for
CSRNet the gap is dramatic ($0.320$ vs.\ Adam's best of $2.74$ at any
learning rate), while for MobileCount it is modest ($0.385$ vs.\
$0.420$). CSRNet's ImageNet-pretrained, batch-norm-free VGG-16
frontend is progressively eroded by Adam's coupled weight decay, which is amplified by the adaptive normalisation
wherever gradients are small, leaving Adam uncompetitive at every
learning rate. MobileCount, which is batch-normalised and trained from
scratch — and therefore largely invariant to the scale of its weights
— tolerates coupled decay far better, and Adam trails AdamW only
modestly. 

The training curves also reveal a behavioural difference: CSRNet's
validation MAE descends to a plateau and remains there, so its final
and best checkpoints nearly coincide, whereas MobileCount reaches its
optimum mid-training and drifts mildly upward thereafter — most
visibly at $10^{-4}$ — so checkpoint selection matters for its
reported accuracy.%maybe put in another section

\begin{table}[t]
    \centering
    \caption{Best validation MAE over the optimiser $\times$
    learning-rate grid ($386\times260$, batch size $6$, $1500$ epochs).
    Asterisks mark the framework defaults.}
    \label{tab:lr_grid}
    \begin{tabular}{ccc@{\hspace{2em}}ccc}
        \toprule
        \multicolumn{3}{c}{CSRNet} & \multicolumn{3}{c}{MobileCount} \\
        \midrule
        LR & Adam* & AdamW & LR & Adam* & AdamW \\
        $10^{-6}$  & $2.74$ & $0.923$          & $10^{-5}$  & $0.880$ & $0.831$ \\
        $10^{-5}$* & $3.51$ & $0.394$          & $10^{-4}$* & $0.632$ & $0.542$ \\
        $10^{-4}$  & $4.33$ & $\mathbf{0.320}$ & $10^{-3}$  & $0.420$ & $\mathbf{0.385}$ \\
        \bottomrule
    \end{tabular}
\end{table}

\begin{figure*}[t]
    \centering
    \includegraphics[width=\linewidth]{images/overfit_3panel_merged.png}
    \caption{Validation MAE over training for the optimiser $\times$
learning-rate grid at $386\times260$ (top: CSRNet; bottom:
MobileCount; columns: learning rates). Thin lines show per-epoch
values, thick lines a running median; stars mark the best checkpoint
of each run, and dotted lines its MAE level.
    $\times$ learning-rate grid at $386\times260$.}
    \label{fig:lrxopt_csrnet}
\end{figure*} %edit the image --make it run to 1500? idk

\paragraph{Batch size.}
\label{sec:exp_batchsize}
After fixing the optimiser and learning rate, batch size was tuned.
Because the optimal learning rate scales with batch size, each batch
size $B\in\{1,2,4,8\}$ was evaluated at its own $\sqrt{B/6}$-scaled
learning rate, anchored to each network's batch-size-$6$
optimum~\cite{hoffer2017train}; this keeps the scale of the per-step
gradient noise approximately constant, so the sweep isolates the
batch-size effect rather than conflating it with an untuned learning
rate.
% TODO: add hoffer2017train (or your preferred sqrt-scaling reference).

The result (Table~\ref{tab:batchsize}) is a clean divergence:
\textbf{CSRNet prefers small batches; MobileCount prefers large ones.}
For CSRNet, MAE improves steadily as the batch shrinks, with the
benefit concentrated below batch size $4$; batch size $1$ ($0.288$) is
the best result obtained at this resolution. MobileCount's optimum
sits at the batch-size-$6$ anchor ($0.385$), with batch sizes $4$ and
$8$ close behind, while batch size $1$ collapses to MAE $2.85$. The
dominant cause is architectural: MobileCount's batch-normalisation
statistics degenerate when each batch contains a single sample.
CSRNet, whose VGG-16 frontend contains no batch normalisation, is
immune to this failure mode — consistent with the original
publication's choice of batch size $1$~\cite{li2018csrnet}.

\begin{table}[h]
\centering
\caption{Batch-size sweep at $\sqrt{B/6}$-scaled learning rates (AdamW,
$386\times260$, $1500$ epochs; best validation MAE). The batch-size-$6$
rows are the anchor cells from Table~\ref{tab:lr_grid}; best per
network in bold.}
\label{tab:batchsize}
\begin{tabular}{l cc@{\hspace{2em}}cc}
\toprule
 & \multicolumn{2}{c}{CSRNet} & \multicolumn{2}{c}{MobileCount} \\
\cmidrule(r{2em}){2-3}\cmidrule{4-5}
Batch size & LR & MAE & LR & MAE \\
\midrule
$1$   & $4.0\times10^{-5}$ & $\mathbf{0.288}$ & $4.0\times10^{-4}$ & $2.85$ \\
$2$   & $6.0\times10^{-5}$ & $0.292$          & $6.0\times10^{-4}$ & $0.543$ \\
$4$   & $8.0\times10^{-5}$ & $0.307$          & $8.0\times10^{-4}$ & $0.436$ \\
$6$*  & $1.0\times10^{-4}$ & $0.320$          & $1.0\times10^{-3}$ & $\mathbf{0.385}$ \\
$8$   & $1.2\times10^{-4}$ & $0.319$          & $1.2\times10^{-3}$ & $0.460$ \\
\bottomrule
\end{tabular}
\end{table}


\subsubsection{Final Selected Configurations}
\label{subsec:final_configs}

The selected configurations for both networks are summarised in
Table~\ref{tab:final_configs}. Relative to the framework baselines
trained under the identical protocol (the asterisked cells of
Table~\ref{tab:lr_grid}), tuning improves the best validation MAE from
$3.51$ to $0.288$ for CSRNet and from $0.632$ to $0.385$ for
MobileCount. Notably, both networks end up on the same optimiser but
retain their architecture-specific learning rate and batch size —
neither network's optimum would have been found by tuning on the other
and transferring.

\begin{table}[h]
    \centering
    \caption{Final selected configurations.}
    \label{tab:final_configs}
    \begin{tabular}{lcc}
        \toprule
        & CSRNet & MobileCount \\
        \midrule
        Optimiser             & \multicolumn{2}{c}{AdamW} \\
        Initial learning rate & $4\times10^{-5}$ & $1\times10^{-3}$ \\
        Batch size            & $1$ & $6$ \\
        Weight decay          & \multicolumn{2}{c}{$1\times10^{-4}$} \\
        LR schedule           & \multicolumn{2}{c}{step $\gamma=0.995$} \\
        Training budget       & \multicolumn{2}{c}{$1500$ epochs} \\
        Augmentation          & \multicolumn{2}{c}{online $80\%$ crop} \\
        Density-target scale $k$ & $25.05$ & $638.7$ \\
        & \multicolumn{2}{c}{\footnotesize at $386\times260$; scaled per Eq.~\eqref{eq:logpara_scaling}} \\
        Random seed           & \multicolumn{2}{c}{$3035$} \\
        \bottomrule
    \end{tabular}
\end{table}

\bibliographystyle{IEEEtran}
\bibliography{references}


\end{document}


